\documentclass[../../main.tex]{subfiles}

\input{0-theme/default}

\begin{document}


%%%%%%%%%%%%%%%%%%%%%%%%%%%%%%%%%%%%%%%%%%%%%%%%%%%%%%%%%%%%%%%%%
\section{Introducción a la  Biomasa}

\begin{frame}{Panorama General de la Biomasa}
\justifying
La búsqueda de fuentes de energía renovable y sostenible ha devuelto protagonismo a la \textbf{biomasa}, relegada desde la Revolución Industrial tras el auge del carbón. Actualmente (2025), la biomasa representa cerca del \textbf{12\% del consumo mundial de energía primaria}, con una tendencia creciente debido a la preocupación por el cambio climático y la sostenibilidad.

\vspace{0.5em}
Los principales procesos de conversión son:
\begin{itemize}
  \item \textbf{Combustión directa}
  \item \textbf{Gasificación}
  \item \textbf{Carbonización y torrefacción}
  \item \textbf{Pirólisis}
  \item \textbf{Fermentación y transesterificación} para biocombustibles líquidos
\end{itemize}
\end{frame}

\begin{frame}{Carbonización y Torrefacción}
\justifying
\textbf{Carbonización:} técnica ancestral usada desde 4000 a.C. para obtener carbón vegetal a partir de biomasa, utilizada para la fundición de hierro en India y China. Aún se usa como combustible sin humo o para filtración.

\vspace{0.5em}
\textbf{Torrefacción:} proceso térmico más reciente ("tostado") que mejora la densidad energética y la manejabilidad de la biomasa sólida.
\begin{itemize}
  \item Produce biocarbón con propiedades similares al carbón mineral.
  \item Ideal para \textbf{co-combustión} con carbón en plantas térmicas.
  \item Potencial uso como sustituto del coque en metalurgia.
\end{itemize}
\end{frame}

\begin{frame}{Gasificación de Biomasa}
\justifying
Proceso químico que convierte materiales carbonosos (biomasa, residuos orgánicos) en \textbf{combustibles gaseosos útiles}, como el \textbf{gas de síntesis (syngas)}.

\vspace{0.5em}
\textbf{Características:}
\begin{itemize}
  \item Requiere atmósfera \textbf{reductora} (bajo oxígeno).
  \item Aporta energía en los enlaces químicos del producto.
  \item A diferencia de la combustión, no libera toda la energía térmica de forma inmediata.
\end{itemize}

\textbf{Aplicaciones:} producción de electricidad, biocombustibles, y materias primas químicas.
\end{frame}

\begin{frame}{Pirólisis de Biomasa}
\justifying
\textbf{Pirólisis:} descomposición térmica de la biomasa en ausencia de oxígeno. Produce:
\begin{itemize}
  \item \textbf{Líquidos:} bio-oil (aceite de pirólisis)
  \item \textbf{Sólidos:} carbón
  \item \textbf{Gases:} combustibles
\end{itemize}

\vspace{0.5em}
\textbf{Importancia:} fue la base para la obtención de keroseno, primer combustible líquido transportable. En 2025, sigue siendo clave para producción descentralizada de energía líquida a partir de residuos.
\end{frame}

\begin{frame}{Avances Recientes en Gasificación y Pirólisis}
\justifying
\textbf{Gasificación avanzada:}
\begin{itemize}
  \item \textbf{Gasificación supercrítica:} permite tratar biomasa húmeda (sin secado previo).
  \item \textbf{Hidrogasificación:} adición de hidrógeno para mejorar la relación H/C del combustible.
  \item Uso de \textbf{residuos pesados de petróleo} para obtener hidrocarburos más ligeros.
\end{itemize}

\textbf{Pirólisis avanzada:}
\begin{itemize}
  \item Mayor eficiencia térmica mediante control preciso de temperatura y tiempo.
  \item Integración con gasificación en \textbf{biorefinerías híbridas}.
\end{itemize}
\end{frame}

\begin{frame}{Importancia Industrial y Ambiental}
\justifying
\textbf{Objetivos de estos procesos:}
\begin{itemize}
  \item No solo conversión energética, sino \textbf{producción de insumos químicos}.
  \item Disminución de residuos.
  \item Reducción de emisiones de gases de efecto invernadero.
\end{itemize}

\textbf{En 2025:} Se están construyendo y ampliando \textbf{plantas industriales de gasificación} para convertir carbón, residuos urbanos y biomasa en fertilizantes, combustibles y plásticos sostenibles.
\end{frame}

\begin{frame}{¿Qué es la Biomasa?}
\justifying
La \textbf{biomasa} proviene de especies vivas, como plantas y animales, o de organismos que estuvieron vivos recientemente. A diferencia de los combustibles fósiles, su formación no requiere millones de años.

\vspace{0.5em}
Las plantas convierten el \textbf{CO\textsubscript{2}} y el agua en compuestos orgánicos mediante la \textbf{fotosíntesis}. Los animales obtienen biomasa al alimentarse de plantas u otros animales.

\vspace{0.5em}
\textbf{Ventaja clave:} La biomasa es \textbf{renovable} y \textbf{reproducible}, lo que la hace atractiva como fuente energética y materia prima química.
\end{frame}

\begin{frame}{Biomasa y Carbono Neutralidad}
\justifying
Cada año se genera gran cantidad de biomasa por fotosíntesis, que absorbe \textbf{CO\textsubscript{2}} atmosférico. Al quemarla, libera ese mismo carbono, por lo que no se incrementa el CO\textsubscript{2} global.

\vspace{1em}
\textbf{Por eso:} la biomasa es considerada una fuente \textbf{carbono-neutral}.

\vspace{1em}
En 2025, la biomasa representa alrededor del \textbf{10–12\%} del consumo primario de energía en el mundo, con un potencial técnico para alcanzar hasta el \textbf{25–30\%} según la IEA.
\end{frame}


\begin{frame}{Tipos de Biomasa y Conversión}
\justifying
La biomasa abarca una gran variedad de materiales:

\begin{itemize}
  \item Desde pastos hasta árboles y residuos agrícolas.
  \item Desde insectos hasta excrementos animales.
\end{itemize}

\vspace{0.5em}
Se distinguen tres tipos principales de productos energéticos obtenidos:

\begin{itemize}
  \item \textbf{Combustibles líquidos:} etanol, biodiésel, metanol, aceite vegetal.
  \item \textbf{Combustibles gaseosos:} biogás, syngas, gas de síntesis.
  \item \textbf{Combustibles sólidos:} carbón vegetal, biocarbón, biomasa torrefactada.
\end{itemize}
\end{frame}

\begin{frame}{Biomasa Alimentaria vs No Alimentaria}
\justifying
\textbf{Ejemplos de biomasa alimentaria:}
\begin{itemize}
  \item Maíz, caña de azúcar, remolacha → etanol
  \item Soya, colza, palma → biodiésel
\end{itemize}

\textbf{Ejemplos de biomasa lignocelulósica:}
\begin{itemize}
  \item Paja, cáscaras, madera, residuos agrícolas y forestales
  \item Utilizados para producir etanol, biogás o combustibles sólidos
\end{itemize}

\vspace{0.5em}
\textbf{En 2025:} se prioriza el uso de residuos no alimentarios para evitar tensiones en la seguridad alimentaria mundial.
\end{frame}

\begin{frame}{Usos Industriales de la Biomasa}
\justifying
\textbf{1. Industria química:}
\begin{itemize}
  \item Producción de metanol, fertilizantes, fibras sintéticas, bioplásticos.
  \item Plataformas basadas en azúcares (glucosa, xilosa) y syngas (CO + H\textsubscript{2}).
\end{itemize}

\textbf{2. Industria energética:}
\begin{itemize}
  \item Generación de calor y electricidad.
  \item Ejemplos: calderas de lecho fluido, cogeneración.
\end{itemize}
\end{frame}

\begin{frame}{Biomasa en el Transporte}
\justifying
\textbf{Combustibles de biomasa para transporte:}
\begin{itemize}
  \item \textbf{Etanol:} obtenido por fermentación de azúcares (caña, maíz).
  \item \textbf{Biodiésel:} extraído y esterificado de aceites vegetales (colza, soya).
  \item \textbf{Bio-oil:} a partir de pirólisis y tratado para obtener “diesel verde”.
\end{itemize}

\vspace{0.5em}
\textbf{También se utiliza:} metano generado por digestión anaerobia o gasificación, como sustituto del gas natural vehicular.
\end{frame}

\begin{frame}{Aplicaciones Ambientales}
\justifying
\textbf{1. Carbón activado:} filtración de agua y eliminación de metales pesados (como mercurio en centrales térmicas).

\textbf{2. Biocarbón (biochar):}
\begin{itemize}
  \item Producido por pirólisis.
  \item Permite el secuestro de carbono en suelos.
  \item Mejora la fertilidad del suelo y aumenta la retención de agua.
\end{itemize}

\vspace{0.5em}
\textbf{Potencial estimado:} más de 0.16 Pg/año de biochar a partir de residuos agrícolas y urbanos (Lehmann et al., 2006).
\end{frame}


%%%%%%%%%%%%%%%%%%%%%%%%%%%%%%%%%%%%%%%%%%%%%%%%%%%%%%%%%%%%%%%%%%%%%%%
\section{Conversión de la Biomasa}
\begin{frame}{Conversión de la Biomasa}
\justifying
La biomasa, en su forma sólida y voluminosa, es difícil de manejar, transportar y usar directamente, especialmente como combustible para el transporte.

\vspace{0.5em}
Por ello, se transforma en \textbf{combustibles líquidos o gaseosos} mediante dos rutas principales:

\begin{itemize}
  \item \textbf{Ruta bioquímica} (fermentación, digestión)
  \item \textbf{Ruta termoquímica} (pirólisis, gasificación, combustión)
\end{itemize}
\end{frame}
\begin{frame}{Conversión de biomasa}
    \begin{figure}
        \centering
        \includegraphics[width=0.5\linewidth]{Energias Alternativas/3-Biomass/img/Conversion_biomasa.png}
        \caption{Conversión de biomasa}
        %\label{fig:enter-label}
    \end{figure}
\end{frame}

\begin{frame}{1. Ruta Bioquímica}
\justifying
Descompone la biomasa mediante microorganismos o enzimas, a baja temperatura y sin necesidad de mucha energía externa.

\vspace{0.5em}
\textbf{Procesos principales:}
\begin{itemize}
  \item \textbf{Digestión anaerobia y aerobia} → genera biogás (\(CH_4\) + \(CO_2\))
  \item \textbf{Fermentación} → produce etanol a partir de azúcares
  \item \textbf{Hidrólisis} (ácida o enzimática) → convierte celulosa en azúcares fermentables
\end{itemize}

\vspace{0.5em}
\textbf{Ventaja:} Tecnología madura para residuos ricos en almidón o azúcar (ej. maíz, caña).
\end{frame}
\begin{frame}{Fermentación e Hidrólisis}
\justifying
\textbf{Fermentación:}
\begin{itemize}
  \item Se convierte el azúcar en etanol usando levaduras.
  \item Funciona bien con cultivos como maíz o caña (ricos en almidón o azúcar).
\end{itemize}

\textbf{Limitaciones:}
\begin{itemize}
  \item No se aprovecha la lignina.
  \item No es eficiente con biomasa lignocelulósica (ej. madera, tallos).
\end{itemize}

\textbf{Hidrólisis:}
\begin{itemize}
  \item Convierte celulosa en azúcares fermentables.
  \item Puede ser \textbf{ácida} (tecnología madura) o \textbf{enzimática} (más eficiente, pero costosa).
\end{itemize}
\end{frame}
\begin{frame}{2. Ruta Termoquímica}
\justifying
Convierte la biomasa completa en gases o líquidos mediante reacciones químicas a alta temperatura.

\vspace{0.5em}
\textbf{Procesos principales:}
\begin{itemize}
  \item \textbf{Combustión:} con exceso de oxígeno → calor y vapor
  \item \textbf{Gasificación:} con poco oxígeno → gas de síntesis (\(CO_2\) + \(H_2\))
  \item \textbf{Pirólisis:} sin oxígeno → bioaceite, gases y carbón vegetal
  \item \textbf{Licuación:} descompone moléculas grandes en líquidos → requiere catalizador
\end{itemize}
\end{frame}
\begin{frame}{Comparación Bioquímica vs Termoquímica}
\justifying
\textbf{Bioquímica:}
\begin{itemize}
  \item Más desarrollada comercialmente para etanol (maíz, caña).
  \item Procesos lentos (por lotes).
  \item No usa toda la biomasa (lignina queda fuera).
\end{itemize}

\textbf{Termoquímica:}
\begin{itemize}
  \item Procesos más rápidos y continuos.
  \item Usa la biomasa completa (incluye lignocelulosa).
  \item Requiere temperaturas elevadas y equipos especiales.
\end{itemize}
\end{frame}
\begin{frame}{Ejemplo: Conversión de una Planta de Maíz}
\justifying
En la conversión bioquímica:

\begin{itemize}
  \item Solo el \textbf{grano} (almidón) se usa para producir etanol.
  \item El \textbf{resto de la planta} (tallos, hojas, raíces) no se aprovecha directamente.
\end{itemize}

En la ruta termoquímica:

\begin{itemize}
  \item Toda la biomasa se puede convertir en \textbf{syngas}.
  \item Luego puede sintetizarse etanol, metanol o combustibles líquidos.
\end{itemize}
\end{frame}
\begin{frame}{Aplicación de la biomasa}
    \begin{figure}
        \centering
        \includegraphics[width=0.55\linewidth]{Energias Alternativas/3-Biomass/img/Aplicacion_Biomasa.png}
        \caption{Aplicaciones de la biomasa}
        %\label{fig:enter-label}
    \end{figure}
\end{frame}
\section{Combustión Directa}
\begin{frame}{¿Qué es la Combustión Directa?}
\justifying
La \textbf{combustión directa} es el método más antiguo para extraer energía de la biomasa. Consiste en la \textbf{reacción química de la materia orgánica con el oxígeno}, generando calor (reacción exotérmica).

\vspace{0.5em}
\textbf{Productos de la combustión:}
\begin{itemize}
  \item Dióxido de carbono (CO\textsubscript{2})
  \item Agua (H\textsubscript{2}O)
  \item Óxidos de azufre y nitrógeno (si están presentes)
\end{itemize}
\end{frame}
\begin{frame}{Factores que Afectan la Combustión}
\justifying
\textbf{1. Características físicas de la biomasa:}
\begin{itemize}
  \item \textbf{Humedad:} se recomienda < 15\%  para una eficiencia (<80\%). Si es > 50\% para una eficiencia (<60\%), reduce la eficiencia térmica.
  \item \textbf{Tamaño de partícula (granulometría)} y \textbf{densidad}
\end{itemize}

\textbf{2. Características químicas:}
\begin{itemize}
  \item Bajo contenido en azufre → bajas emisiones contaminantes.
\end{itemize}

\textbf{3. Condiciones de operación:}
\begin{itemize}
  \item Temperatura: entre 600 °C y 1300 °C.
  \item Exceso de oxígeno para evitar formación de monóxido de carbono (CO).
\end{itemize}
\end{frame}



\begin{frame}{Generación de vapor por combustión directa}
    \begin{figure}
        \centering
        \includegraphics[width=0.55\linewidth]{Energias Alternativas/3-Biomass/img/Generacion_combustion_directa.png}
        \caption{Generación de vapor por combustión directa}
        %\label{fig:enter-label}
    \end{figure}
\end{frame}
\section{Procesos Termoquímicos}
\begin{frame}{Procesos Termoquímicos}
\justifying
La biomasa puede convertirse en combustibles sólidos, líquidos o gaseosos mediante procesos termoquímicos, que implican:

\begin{itemize}
  \item Altas temperaturas (260 °C – 1500 °C)
  \item Condiciones de oxidación controlada
  \item Diferentes presiones (1 a 30 atm)
\end{itemize}

\vspace{0.5em}
Los principales procesos son:
\begin{enumerate}
  \item \textbf{Gasificación}
  \item \textbf{Pirólisis}
\end{enumerate}
\end{frame}
\begin{frame}{Gasificación}
\justifying
La \textbf{gasificación} es la oxidación parcial de biomasa sólida para obtener un gas combustible compuesto por:

\begin{itemize}
  \item CO(monóxido de carbono), CO\textsubscript{2}(dióxido de carbono), H\textsubscript{2}(hidrógeno), CH\textsubscript{4}(metano), y vapor de agua
\end{itemize}

\vspace{0.5em}
\textbf{Condiciones óptimas:}
\begin{itemize}
  \item Temperatura: 700°C – 1500°C
  \item Presión: 1 – 30 atm
  \item Humedad baja (< 15\%) para mayor eficiencia
\end{itemize}
\end{frame}


\begin{frame}{Tipos de Gasificadores}
\justifying
\begin{columns}
  % Columna izquierda: texto
  \begin{column}{0.55\textwidth}
    Los gasificadores se clasifican en:

    \begin{itemize}
      \item \textbf{Lecho fijo:}
        \begin{itemize}
          \item Flujo ascendente
          \item Flujo descendente
        \end{itemize}
      \item \textbf{Lecho fluidizado:}
        \begin{itemize}
          \item Mejor mezcla de aire y biomasa
          \item Más eficiente y flexible
        \end{itemize}
    \end{itemize}

    \vspace{0.5em}
    El gas generado puede usarse directamente en turbinas o transformarse en combustibles líquidos como metanol.
  \end{column}

  % Columna derecha: imagen con entorno figure
  \begin{column}{0.45\textwidth}
    \begin{figure}
      \centering
      \includegraphics[width=\linewidth]{Energias Alternativas/3-Biomass/img/Gasificadores.png}
      \caption{Tipos de gasificadores}
    \end{figure}
  \end{column}
\end{columns}
\end{frame}


\begin{frame}{Tipos de Gas de Gasificación}
\justifying
\textbf{Gas gasógeno (con aire):}
\begin{itemize}
  \item Bajo poder calorífico (< 5.5 MJ/m³)
  \item Alta proporción de nitrógeno
  \item Uso directo en turbinas para generar electricidad
\end{itemize}

\textbf{Gas de síntesis (con oxígeno):}
\begin{itemize}
  \item Poder calorífico medio
  \item Rico en CO y H₂
  \item Puede transformarse en combustibles líquidos como metanol
\end{itemize}
\end{frame}
\begin{frame}{Pirólisis}
\justifying
La \textbf{pirólisis} es la descomposición térmica de la biomasa en ausencia total de oxígeno.

\textbf{Condiciones:}
\begin{itemize}
  \item Temperatura: 260 °C – 550 °C (convencional)
  \item Hasta 1000 °C en residuos sólidos urbanos (RSU)
\end{itemize}

\textbf{Productos:}
\begin{itemize}
  \item Carbón vegetal
  \item Combustible líquido (bio-oil)
  \item Gas combustible
\end{itemize}
\end{frame}
\begin{frame}{Pirólisis Súbita (Fast)}
\justifying
Técnica avanzada de pirólisis a alta temperatura y tiempo corto.

\textbf{Características:}
\begin{itemize}
  \item Temperatura: 800 °C – 1000 °C
  \item Duración: segundos
\end{itemize}

\textbf{Resultados:}
\begin{itemize}
  \item 60\% gas rico en H₂ y CO
  \item 10\% sólidos
  \item Resto: líquidos (bio-oil)
\end{itemize}

\vspace{0.5em}
Ideal para el tratamiento de \textbf{residuos sólidos urbanos}.
\end{frame}





%%%%%%%%%%%%%%%%%%%%%%%%%%%%%%%%%%%%%%%%%%%%%%%%%%%%%%%%%%%%%%%%%%

\section{Procesos Biológicos de Conversión}
\begin{frame}{Procesos Biológicos de Conversión}
\justifying
Los procesos biológicos convierten biomasa húmeda en energía gracias a la acción de \textbf{microorganismos}, sin requerir altas temperaturas.

\vspace{0.5em}
\textbf{Procesos principales:}
\begin{itemize}
  \item \textbf{Digestión anaerobia} → biogás (CH\textsubscript{4}metano + CO\textsubscript{2}dióxido de carbono)
  \item \textbf{Fermentación alcohólica} → etanol o metanol
\end{itemize}

\vspace{0.5em}
Estos procesos producen combustibles renovables con alta densidad energética.
\end{frame}
\begin{frame}{Digestión Anaerobia}
\justifying
Proceso biológico que descompone biomasa en ausencia de oxígeno, mediante bacterias.

\vspace{0.5em}
\textbf{Etapas del proceso:}
\begin{itemize}
  \item \textbf{Hidrólisis:} macromoléculas → azúcares
  \item \textbf{Acetogénesis:} azúcares → ácidos orgánicos
  \item \textbf{Metanogénesis:} ácidos → CH\textsubscript{4}, CO\textsubscript{2}, SH\textsubscript{2} (sulfuro de hidrógeno)
\end{itemize}

\vspace{0.5em}
El producto es \textbf{biogás}, y un residuo sólido (fango) con valor como fertilizante o alimento animal.
\end{frame}
\begin{frame}{Esquema Conceptual del proceso de digestión anaeróbica}
    \begin{figure}
        \centering
        \includegraphics[width=0.55\linewidth]{Energias Alternativas/3-Biomass/img/digestion_anaerobica.png}
        \caption{proceso de digestión anaeróbica}
        \label{fig:1}
    \end{figure}
\end{frame}

\begin{frame}{Tipos de Digestores Anaerobios}
\justifying
\textbf{Digestores discontinuos:}
\begin{itemize}
  \item Proceso por lotes (ya en desuso)
  \item Tiempo de residencia: 20 a 60 días
\end{itemize}

\textbf{Digestores continuos:}
\begin{itemize}
  \item Operan sin interrupción
  \item Tipos:
    \begin{itemize}
      \item Mezcla completa (TRH: 15–25 días)
      \item Contacto (TRH: 4–8 días)
      \item Filtro anaerobio (TRH: 1–4 días)
    \end{itemize}
\end{itemize}

\vspace{0.5em}
\textbf{TRH} = Tiempo de Residencia Hidráulico.
\end{frame}
\begin{frame}{Tipos de digestores}
    \begin{figure}
        \centering
        \includegraphics[width=0.5\linewidth]{Energias Alternativas/3-Biomass/img/tipos_digestores.png}
        \caption{Tipos de digestores}
        %\label{fig:enter-label}
    \end{figure}
\end{frame}
\begin{frame}{Aplicaciones del Biogás}
\justifying
El biogás (CH\textsubscript{4} + CO\textsubscript{2}) generado por digestión anaerobia se usa como:

\begin{itemize}
  \item \textbf{Fuente de calor} (cocción, calefacción)
  \item \textbf{Combustible para calderas} (generación de vapor)
  \item \textbf{Combustible en motores} (generación de electricidad)
\end{itemize}

\vspace{0.5em}
El fango residual puede aprovecharse como \textbf{fertilizante} o incluso en \textbf{alimentación animal}.
\end{frame}
\begin{frame}{Fermentación Alcohólica}
\justifying
Transformación de azúcares simples en alcohol por acción de microorganismos (levaduras).

\vspace{0.5em}
\textbf{Biocombustibles obtenidos:}
\begin{itemize}
  \item \textbf{Etanol} (CH\textsubscript{3}-CH\textsubscript{2}-OH)( alcohol etílico).
  \item \textbf{Metanol} (CH\textsubscript{3}OH) (alcohol metílico)
\end{itemize}

\vspace{0.5em}
Usado como combustible en:
\begin{itemize}
  \item \textbf{Motores a gasolina} (uso directo o mezclado)
  \item \textbf{Motores diésel} (requiere modificaciones)
\end{itemize}
\end{frame}
\begin{frame}{Tratamiento Previo y Hidrólisis}
\justifying
Para convertir carbohidratos complejos (almidón, celulosa) en alcohol se requiere:

\begin{itemize}
  \item \textbf{Pretratamiento mecánico:} trituración, molienda (↑ superficie de contacto)
  \item \textbf{Hidrólisis:}
    \begin{itemize}
      \item \textbf{Enzimática} (con microorganismos)
      \item \textbf{Química} (con ácidos)
    \end{itemize}
\end{itemize}

\vspace{0.5em}
El etanol se destila para eliminar el agua y obtener combustible de alta pureza.
\end{frame}
\begin{frame}{Esquema conceptual del proceso de fermentación alcohólica.}
    \begin{figure}
        \centering
        \includegraphics[width=0.55\linewidth]{Energias Alternativas/3-Biomass/img/fermentación_alcohólica..png}
        \caption{fermentación alcohólica}
        %\label{fig:enter-label}
    \end{figure}
\end{frame}
\section{Residuos Sólidos Urbanos (RSU)}
\begin{frame}{Residuos Sólidos Urbanos (RSU)}
\justifying
Los RSU son residuos generados principalmente por la actividad humana en hogares, y también por pequeñas industrias.

\vspace{0.5em}
\textbf{Problema:} El aumento del consumo → incremento de residuos → amenaza ambiental.

\vspace{0.5em}
\textbf{Composición típica:}
\begin{itemize}
  \item 50\% materia orgánica (fermentable)
  \item 35\% materiales combustibles (papel, madera, plásticos)
  \item 15\% vidrio, metales, etc.
\end{itemize}
\end{frame}
\begin{frame}{Opciones para la Eliminación de RSU}
\justifying
\textbf{Sin recuperación:}
\begin{itemize}
  \item Vertedero sin aprovechamiento
\end{itemize}

\textbf{Con recuperación:}
\begin{itemize}
  \item \textbf{Aprovechamiento energético:}
    \begin{itemize}
      \item Incineración
      \item Digestión anaerobia
      \item Gases de vertedero
      \item Pirólisis
    \end{itemize}
  \item \textbf{Aprovechamiento de materiales:}
    \begin{itemize}
      \item Reciclado (papel, vidrio, metal, plásticos)
      \item Producción de fertilizantes
    \end{itemize}
\end{itemize}
\end{frame}
\begin{frame}{Incineración de Residuos}
\justifying
\textbf{Proceso:}
\begin{enumerate}
  \item Los residuos son transportados y almacenados.
  \item Se incineran en hornos (parrilla, rotativos, lecho fluido).
  \item El calor se transfiere a una caldera → genera vapor.
  \item El vapor acciona una turbina que produce electricidad.
\end{enumerate}

\vspace{0.5em}
\textbf{Ventajas:} Reducción de volumen, generación de energía.
\textbf{Requiere:} Sistemas de control de emisiones contaminantes.
\end{frame}
\begin{frame}{Tratamiento de Gases Contaminantes}
\justifying
La combustión de RSU genera gases tóxicos que deben ser tratados antes de su emisión.

\vspace{0.5em}
\textbf{Contaminantes principales:}
\begin{itemize}
  \item Dioxinas y furanos
\end{itemize}

\vspace{0.5em}
\textbf{Métodos de tratamiento:}
\begin{itemize}
  \item Pulverización con cal + agua (semihúmedo)
  \item Lavado y secado de gases
\end{itemize}

\vspace{0.5em}
\textbf{Residuos finales:} Escorias y cenizas → se trasladan a vertederos de seguridad.
\end{frame}
\begin{frame}{Captación de Gases de Vertedero}
\justifying
La materia orgánica enterrada en vertederos se descompone y libera \textbf{gas metano (CH\textsubscript{4})}.

\vspace{0.5em}
\textbf{Captación:}
\begin{itemize}
  \item Pozos perforados verticales
  \item Tubos laterales conectados a soplantes
\end{itemize}

\vspace{0.5em}
\textbf{Uso:} El metano captado puede convertirse en electricidad mediante sistemas modulares de generación (ej. Highland Energy Inc.).
\end{frame}
\section{Factores que Influyen en los Costes de la Energía de la Biomasa}
\begin{frame}{Factores que Influyen en los Costes de la Energía de la Biomasa}
\justifying
La rentabilidad depende de:
\begin{itemize}
  \item \textbf{Tipo de biomasa} (residuos, cultivos energéticos, etc.)
  \item \textbf{Tecnología empleada} (incineración, digestión, conversión, etc.)
  \item \textbf{Externalidades:} costes asociados a daños a la salud y medioambientales.
\end{itemize}

\vspace{0.5em}
En aplicaciones domésticas, el coste varía según:
\begin{itemize}
  \item Uso (cocina, estufa, chimenea)
  \item Potencia térmica instalada
  \item Tipo y coste del combustible (ej. briquetas, pellets: 0,15–0,30 €/kg)
\end{itemize}
\end{frame}
\begin{frame}{Costes en Aplicaciones Colectivas y Producción Eléctrica}
\justifying
\textbf{Aplicaciones colectivas:}
\begin{itemize}
  \item Costes de inversión en la planta térmica: dependen de la capacidad instalada y el número de viviendas.
  \item Ejemplo: Para 1.000 viviendas en clima suave, inversión promedio ≈ 700 €/vivienda.
  \item Red de distribución: inversión ≈ 70 €/metro (para 10 km de red).
\end{itemize}

\vspace{0.5em}
\textbf{Plantas de producción eléctrica:}
\begin{itemize}
  \item El coste del kWh depende de la amortización de la inversión, combustible y O\&M.
  \item El coste del combustible (biomasa) representa aproximadamente el 50\% del coste de la unidad energética.
  \item La biomasa se compra en "termias" (1 termia = 1.000 kcal), a ≈ 0,01 €/termia.
\end{itemize}
\end{frame}
\begin{frame}{Inversión en Instalaciones de Gran Escala}
\justifying
Para una planta de producción eléctrica con biomasa:
\begin{itemize}
  \item Capacidad instalada de 40 MW → inversión aproximada de 50 millones de euros.
  \item Esta inversión puede representar alrededor del 30\% del coste de la unidad energética generada.
\end{itemize}

\vspace{0.5em}
En la producción de biocombustibles:
\begin{itemize}
  \item Bioetanol: 
    \begin{itemize}
      \item Planta de 100.000 m³/año: inversión ≈ 55.000 €
      \item Planta de 20.000 m³/año: inversión superior a 12.000 €
    \end{itemize}
  \item Biodiésel:
    \begin{itemize}
      \item Planta de 30.000 m³/año: inversión ≈ 2 millones de euros
      \item Planta de 5.000 m³/año: inversión ≈ 250.000 €
    \end{itemize}
\end{itemize}
\end{frame}
\begin{frame}{Costes en la Producción de Biogás}
\justifying
El biogás se obtiene mediante la digestión anaerobia de residuos (ganaderos y RSU).  
Para su producción:
\begin{itemize}
  \item Solo se consideran los costes de acondicionamiento de la instalación (inferiores al 25\% del coste global).
  \item Se optimiza el aprovechamiento de la materia orgánica, reduciendo el coste global.
\end{itemize}

\vspace{0.5em}
Estos sistemas permiten convertir residuos en energía, ayudando a paliar problemas medioambientales.
\end{frame}
\section{Ventajas del Uso de la Biomasa}
\begin{frame}{Ventajas del Uso de la Biomasa}
\justifying
\textbf{Beneficios ambientales:}
\begin{itemize}
  \item Reducción de gases de efecto invernadero (hasta 65\%)
  \item Disminución de la contaminación atmosférica
  \item Mejora en la eficiencia energética
  \item Biodegradables y menos contaminantes en caso de vertidos
  \item Menor generación de emisiones como SO₂ (lluvia ácida)
\end{itemize}

\textbf{Beneficios técnicos y económicos:}
\begin{itemize}
  \item Mayor vida útil del motor (actúa como lubricante)
  \item Producción local y diversificación energética
  \item Menor dependencia del petróleo
  \item Transporte y almacenamiento más seguros
\end{itemize}
\end{frame}
\begin{frame}{Criterios de Sostenibilidad para Biocarburantes}
\justifying
\textbf{Directiva Europea sobre Energías Renovables:}
\begin{itemize}
  \item Reducción mínima del 35\% de GEI comparado con combustibles fósiles
  \item Prohibido el uso de:
    \begin{itemize}
      \item Tierras protegidas o con alta biodiversidad
      \item Humedales o zonas con alto contenido de carbono
    \end{itemize}
  \item Cultivo bajo criterios estrictos de sostenibilidad
\end{itemize}

\vspace{0.5em}
\textbf{Acuerdos internacionales:}
\begin{itemize}
  \item En desarrollo para asegurar trazabilidad y sostenibilidad global
\end{itemize}
\end{frame}
\begin{frame}{Otros Impactos Positivos del Uso de Biocombustibles}
\justifying
\begin{itemize}
  \item Alternativa frente al agotamiento de combustibles fósiles
  \item Menor irritabilidad para la piel humana y el medio acuático
  \item Reducción del humo visible en arranque (~30\%)
  \item Disminución de hidrocarburos aromáticos y plomo en combustibles
  \item Aumento del rendimiento agrícola en zonas abandonadas
\end{itemize}
\end{frame}
\begin{frame}{Inconvenientes del Uso de Biocombustibles}
\justifying
\begin{itemize}
  \item Uso de cultivos alimentarios (maíz, caña) para energía
  \item Cambio de uso del suelo (deforestación, pérdida de biodiversidad)
  \item Emisiones indirectas por insumos agrícolas (fertilizantes, transporte)
  \item Alto consumo de agua: hasta 20 kg por cada kg de biodiésel
  \item Puede no ser el mejor uso para ciertos residuos o tierras agrícolas
\end{itemize}
\end{frame}
\begin{frame}{Consideraciones sobre Balance Energético}
\justifying
\textbf{En procesos industriales:}
\begin{itemize}
  \item Emisiones netas no siempre nulas (uso de maquinaria, fertilizantes, transporte)
  \item Energía incorporada en el agua utilizada puede superar la energía generada
\end{itemize}

\vspace{0.5em}
\textbf{A tener en cuenta:}
\begin{itemize}
  \item ¿Es la producción de combustible el mejor uso para ese residuo?
  \item ¿Es ese cultivo el mejor uso del suelo disponible?
\end{itemize}

\vspace{0.5em}
\textit{Cada decisión depende del análisis territorial, social y ambiental específico.}
\end{frame}



\end{document}